% Two-checkpoint tree for Paper B.  The prefix to checkpoint 1 is shared once
% per root; verify-now and continue fork there; the continue branch forks again
% at checkpoint 2 into verify-at-2 and continue.  At most one verifier call per
% root.  No curve in this figure is data.
\begin{figure*}[t]
\centering
\resizebox{\textwidth}{!}{%
\begin{tikzpicture}[
  node distance=5mm and 9mm,
  box/.style={draw, rounded corners, align=center, inner sep=3pt, font=\small, minimum height=6.5mm},
  rnd/.style={box, draw=violet!70!black, fill=violet!8},
  fixed/.style={box, draw=black!60, fill=black!3},
  arrow/.style={-{Latex[length=2mm]}, thick}
]
\node[fixed] (prefix) {shared prefix to\\checkpoint 1 ($\rho$ frozen)};
\node[rnd, right=of prefix] (cp1) {checkpoint-1 capture\\isolated siblings};
\node[fixed, above right=6mm and 12mm of cp1] (now) {$g_{\mathrm{now}}$: verify now\\$\to$ frozen response $\to\rho$};
\node[fixed, below right=6mm and 12mm of cp1] (later) {$g_{\mathrm{later}}$: continue to\\checkpoint 2};
\node[rnd, right=of later] (cp2) {checkpoint-2 capture\\when reached};
\node[fixed, above right=4mm and 9mm of cp2] (v2) {verify at 2\\$\to$ frozen response};
\node[fixed, below right=4mm and 9mm of cp2] (never) {$g_0$: never verify\\$\to\rho$};
\draw[arrow] (prefix) -- (cp1);
\draw[arrow] (cp1) -- (now);
\draw[arrow] (cp1) -- (later);
\draw[arrow] (later) -- (cp2);
\draw[arrow] (cp2) -- (v2);
\draw[arrow] (cp2) -- (never);
\end{tikzpicture}%
}
\par\smallskip
\parbox{0.95\textwidth}{\centering\footnotesize
   At most one verifier call per root; the same frozen verifier and response rule
   are used at checkpoint 1 and checkpoint 2.  Roots that never reach
   checkpoint 2, or at which the verifier is unavailable, stay in the
   original-root analysis under the registered fallback.}
\caption{The two-checkpoint regime tree on one original root.  The prefix to
checkpoint 1 is shared; verify-now and continue fork there; the continue branch
forks again at checkpoint 2 into verify-at-2 and continue.  All original roots,
including early-terminal and infeasible ones, enter the estimand.}
\label{fig:b-architecture}
\end{figure*}
